\section{FINITE DIMENSIONAL MODULES OVER A SPLIT SEMI-SIMPLE LIE ALGEBRA}

In this paragraph, we retain the general notations of §6. We denote by P (resp. Q) the group of weights of R (resp. radical weights of R). We denote by \(P_{+}\) (resp. \(Q_{+}\) ) the set of elements of P (resp. Q) that are positive for the order relation defined by B. We denote by \(P_{++}\) the set of dominant weights of R relative to B (Chap. VI, §1, no. 10). An element \(\lambda\) of \(h^{*}\) belongs to P (resp. to \(P_{++}\) ) if and only if all the \(\lambda(H_{\alpha})\) , \(\alpha \in B\) , are integers (resp. integers \(\geq 0\) ). We have \(P_{++} \subset P_{+}\) (Chap. VI, §1, no. 6). If \(w \in W\) , we denote by \(\varepsilon(w)\) the determinant of w, which is equal to 1 or -1. We put \(\rho = \frac{1}{2} \sum_{\alpha \in R_{+}} \alpha\) .
% semantic-fragments: legacy-input-seam
\relax{}
